% ---------------------------------------------------------------------------
% Heuristic Learning
%
% RENAMED from "Methodology" on 2026-08-19 and split into two subsections at the
% user's direction: Motivation, then Introducing HL. The section now carries the
% DEFINITION, not just the loop.
%
% THE DEFINITION, in the user's own words (2026-08-19): heuristic learning is the
% paradigm of using an agent to understand feedback and directly modify a program
% or other interpretable system. Two things follow that the rest of the paper does
% not yet reflect:
%   - it is NOT game-specific. "Program or other interpretable system" is the
%     general object; adversarial games are the setting we measure it in, not the
%     definition's scope. The introduction currently frames HL game-specifically
%     ("a coding agent maintains a heuristic program ... rescores it against a
%     frozen opponent pool"), which is the instantiation. State the general form
%     here first, then narrow to the game instantiation explicitly, so a reader
%     cannot mistake the benchmark for the boundary of the paradigm.
%   - "understand feedback" and "directly modify" are the two halves, and they map
%     onto the two bounds the paper already argues: what the loop can READ (lower
%     bound, information per sample) and what it can EDIT (upper bound,
%     maintenance capacity). The definition is where those two are introduced as
%     one object; do not let the subsections split them apart.
%
% WRITTEN 2026-08-19: 3.1 (four paragraphs + Figure 1) and 3.2's opening (the
% definition, the two-agent loop, Figure 2, the two costly constraints). What is
% written is the user's own outline: expressiveness and interpretability, the size
% figure, the program space as a convenient destination for human experience, the
% maintenance cost, then the agent-era reversal (input and output space IS program
% text, so both halves of the loop run in one representation), then the loop itself.
%
% TWO JUDGEMENT CALLS in 3.1, both reversible:
%   - The user's phrasing was that HL is "a robust architecture with very high
%     information flow". The written version claims the information flow (that is
%     architectural and follows from the representation argument) but NOT the
%     robustness, which is \val{robustspread} and unmeasured. Calling the
%     architecture robust in the motivation is invariant 9 exactly -- the framing
%     ahead of the number -- and it would be the second time robustness got asserted
%     before it had a result behind it. When \val{robustspread} exists, this is where
%     the claim can be made if it holds.
%   - The maintenance paragraph was rewritten once because its first version
%     restated 01-introduction.tex's "every new case has to be noticed, encoded, and
%     reconciled ... at a cost that grows with the program" almost clause for clause.
%     What this section adds and the intro does not have: the cost is in the
%     CHECKING, not the writing, which is why it grows rather than staying flat. If
%     that distinction is ever cut, cut the paragraph -- without it the paragraph is
%     the intro again.
%
% STILL OWED BY THIS SECTION, none of it written yet: everything in the blocks
% below. The lower-bound argument, the on/off-policy spectrum (offpolicyvaluellm2026,
% samplereuse2022 -- still cited nowhere in body text), the echo2025hindsight
% concession in its empirical form, the description-length protocol, the baseline
% rationale, the budget-commensurability conversion, and the measurement definitions.
% 3.2 currently stops after the loop; the protocol half of it is absent.
%
% MOTIVATION SUBSECTION. The introduction is capped at 2.5 pages and several
% argued links were deleted from it wholesale rather than compressed; the ones
% about WHY the paradigm is worth revisiting land here. See the intro's
% "WHERE THE DELETED MATERIAL WENT" block for the list -- this subsection is the
% first place with room to argue them at length rather than assert them.
%
% The HL loop, one coding_agent_act per policy version: read rules + own source
% + replays of losses -> hypothesize failure modes -> revise policy -> rescore
% on the frozen set. Append-only event log; missing token/time data recorded as
% unknown/missing, never zero.
%
% Also here:
%   - the RL baselines (model-based, policy gradient) and how their budget is
%     made commensurable with HL's. This is the crux of the sample-efficiency
%     claim: an episode and a coding_agent_act are not the same unit, so state
%     the conversion and its assumptions explicitly.
%   - measurement: Final Elo, Elo gain, Elo-budget curve, Elo-budget AUC. Every
%     AUC carries its x-axis unit (AUC_Elo_episode, AUC_Elo_token, ...).
%   - LocalPolicyKL between adjacent versions, over decision points actually
%     reached. Deliberately not framed as epistemic uncertainty reduction.
%     Occupancy shift is logged separately and never summed with the KL.
%   - the discipline: a policy change is not a result. The question is how much
%     Elo a given behaviour change bought.
%
% The description-length protocol lives here, and it has to be stated before any
% result uses it, because the introduction's upper bound rests on it.
% K(pi_good) is uncomputable; what is measured is \val{kproxy}. Fix and state:
%   - compressed bytes, not lines. Lines are a style artifact and the archive
%     spans several languages. Name the compressor and its settings.
%   - comments and dead code stripped before compression, with the stripping
%     tool named. A program's comments are addressed to a human reader and are
%     not part of the policy.
%   - cross-language normalization. Without it, a verbose language looks like a
%     complex policy. State the normalization and show the result survives it.
%   - program size logged at *every* revision, not only at the plateau. The
%     upper-bound claim is about a trajectory approaching a capacity, so the
%     trajectory is the evidence; the endpoint alone cannot distinguish a
%     plateau from a stall.
% Anything measured this way is a proxy and must be called one in the text every
% time the axis is used. Cyclomatic complexity (mccabe1976complexity) is a
% secondary axis at best -- report it if it agrees, do not lead with it.
%
% THE WHOLE LOWER-BOUND ARGUMENT MOVED HERE on 2026-08-18 (second cut, to reach
% 2.5 pages). The introduction now asserts only that HL is more sample-efficient
% at matched Elo; the reason is gone from it entirely and has to be made here,
% attached to the budget-commensurability protocol above, because that is where a
% reader asks why the two units can be compared at all. What this section owes:
%   - information per sample is the binding constraint. A scalar reward reports
%     that an episode went badly; the rules plus a replay of the loss report what
%     happened, when the position turned, and which decision was responsible. The
%     difference compounds, because it sets how much of a sample survives to the
%     next revision.
%   - the on-policy/off-policy spectrum, with offpolicyvaluellm2026 (policy
%     gradient uses a batch once and discards it) and samplereuse2022 (off-policy
%     reuse). Both keys are cited in 04-benchmark.tex's scope comment only, so
%     the argument's two supports currently appear nowhere in the body.
%   - a maintained program sits at the far end: off-policy AND cross-policy-class
%     -- it learns from replays of any opponent and from human source code no
%     gradient method can consume, and it reads the rules directly, a prior
%     costing no samples at all. This is the sentence the baseline rationale below
%     was written to follow.
%   - the concession that keeps the claim honest, and it must not be dropped: the
%     information-per-sample gain is NOT the symbolic form's doing. echo2025hindsight
%     rewrites its own failed episodes into experience and gets the same kind of
%     gain holding no program at all, so what buys it is the reader. The program
%     earns its place on the other side of the question, as what makes the policy
%     editable. 02-preliminary.tex's scope comment already carries the definitional
%     half of this (lower bound over what the loop can READ, upper over what it can
%     EDIT); this section owes the empirical half. Merging the two is what produces
%     the unsupportable version of the claim.
%
% 3.3 BOTTLENECKS, written 2026-08-19 to the user's outline: "observation轨迹（决定
% 了反馈信息）/ revise规则（决定了迭代的鲁棒性和稳定性）/ 大致说这两点". Two points,
% broadly, and the subsection is deliberately short.
%
% WHY THIS AXIS AND NOT ANOTHER: the two bottlenecks the user named are exactly the
% two bounds the paper already argues, seen from the implementation side. The
% observation trace is what the loop can READ (lower bound, information per sample,
% Play/tau_t in Algorithm 1); the revise step is what it can EDIT (upper bound,
% maintenance capacity, Revise). So 3.3 is written as the two load-bearing lines of
% Algorithm 1 and nothing else -- every other line is bookkeeping. Do not let this
% subsection drift into a general limitations list; the limitations of the STUDY are
% 06-discussion.tex, and these are constraints of the PARADIGM.
%
% INVARIANT 9 APPLIES HERE, and this is the second time in this section. The user
% attached "鲁棒性和稳定性" to the revise half, and \val{robustspread} is registered
% but unmeasured. What is written: the MECHANISM by which the revise rule decides
% whether a trajectory is monotone (scope, acceptance, memory of rejects). What is
% NOT written: any claim that our loop is stable, any magnitude, any comparison.
% The last sentence of the revise paragraph points at 05 on purpose. When
% \val{robustspread} exists, the result attaches there and not here.
%
% THE TRADE-OFF PARAGRAPH IS THE POINT OF PUTTING BOTH IN ONE SUBSECTION. Written
% because two separate bottleneck paragraphs would leave a reader thinking they can
% be fixed independently, and they cannot: the discipline that stabilizes the loop
% (rescore, accept or reject) compresses a whole revision's worth of evidence into
% one number, which is the same scalar-channel problem 3.1 argues RL has. If this
% paragraph is ever cut, the subsection should be split in two, because the reason
% they are together is gone.
%
% evotrace2026 IS CITED HERE AND WAS ORPHANED UNTIL NOW (registered in refs.bib,
% cited in no body text as of 2026-08-19). What is used is the one hard number in
% its abstract: ~30% of code lines added during evolutionary search are
% byte-identical re-introductions of previously deleted lines, in nearly every run.
% Checked against the abstract on 2026-08-19 -- what is NOT in the abstract, and so
% must not be attributed to them: the fraction of edits that are net improvements,
% whether edits are local or global, and diversity collapse. Their setting is
% evolutionary search with a population, not our single maintained artifact, so the
% claim made here is that the failure mode EXISTS and is cheap to induce, not that
% our loop exhibits it at that rate.
%
% ONE THING THIS SUBSECTION OWES 04-benchmark.tex, recorded so it does not fall
% through: the trace paragraph says the contents of tau_t are a design variable and
% does not say what OUR tau_t contains. That belongs with the apparatus, and 04's
% scope comment does not currently list it. If 04 is written without specifying the
% trace contents, this subsection is left arguing that a choice matters while the
% paper never states which choice it made.
%
% BASELINE RATIONALE MOVED HERE FROM SECTION 1 on 2026-08-18 under the 2.5-page
% budget. The intro now says only "model-based RL is our primary baseline"; the
% reason has to be stated here in full, because without it the baseline choice
% looks like it was picked to flatter us. Measuring HL against policy gradient
% alone confounds symbolic-versus-neural with on-policy-versus-off-policy: a
% maintained program is off-policy AND cross-policy-class, so beating an on-policy
% method would leave both explanations open. Model-based RL is off-policy, which
% isolates the representational question; policy gradient stays in as the
% on-policy reference point rather than being dropped. Also owed here: the honest
% form of kuang2025generative's efficiency claim, which the intro now calls
% "directional" without saying why -- one family of tasks, no shared budget
% accounting, no adversary. See 07-related-work.tex, which makes the same point
% about their normalization protocol.
% ---------------------------------------------------------------------------

\section{Adversarial Heuristic Learning (AHL)}
\label{sec:hl}

\subsection{Motivation}
\label{sec:hl-motivation}

\begin{CJK*}{UTF8}{gbsn}

符号式范式的吸引力首先来自它的可读、可解释与可编辑性：程序直接写出决策边界，一条规则或一个搜索截断即可表达复杂行为，并能被人检查和局部修改。图~\ref{fig:size}将这种紧凑性与同等强度的学习策略进行对比；这不是说越小越好，而是说明两者的描述长度可能相差数个数量级。

\begin{center}
\fcolorbox{brandaccent}{white}{\parbox{0.93\linewidth}{\small
\textbf{观察 1：可解释且可编辑。} 符号策略把决策写成可检查、可局部修改的程序；每条经验都能直接落到策略实体上，并立即影响后续决策。}}
\end{center}

% PLACEMENT: [b], not [t], and the reason is worth keeping. With [t] this float
% went to the top of the page its reference falls on, which put it above Section 3's
% own heading and inside the tail of Section 2 -- a reader met the figure while
% still reading about K(pi_good). [b] keeps it below the heading. Recheck this
% whenever text above it grows or shrinks; float placement is a function of the page
% break, so any edit in Sections 2 or 3 can move it.
\begin{figure}[b]
\centering
% PLACEHOLDER, drawn 2026-08-19 at the user's request to leave a slot. This is a
% schematic and must be replaced by a measured plot before submission -- it asserts
% a magnitude gap with no data behind it, which is exactly what the numbers guard
% exists to prevent, so the two values are \val keys and will fail the
% camera-ready build until measured. When the real figure lands, keep the axis
% (description length, log scale) and keep both points annotated with the Elo band
% they were measured at; a size comparison without a strength condition is
% meaningless. Consider adding the archive's lower envelope (\val{envslope}) as a
% third element, since it is the same axis measured on humans.
\begin{tikzpicture}[x=1cm,y=1cm]
  \draw[-{Stealth[length=2mm]}] (0,0) -- (10,0);
  \node[anchor=north east,font=\footnotesize] at (10,-0.55)
    {description length (log scale)};
  \node[draw,rounded corners,inner sep=3pt,fill=black!5] (p) at (2,1)
    {\footnotesize symbolic policy};
  \node[draw,rounded corners,inner sep=3pt,fill=black!5] (n) at (7.6,1)
    {\footnotesize learned policy};
  \draw[dashed] (p) -- (2,0.15);   \draw[dashed] (n) -- (7.6,0.15);
  \node[below] at (2,0)   {\footnotesize \val{progsize}};
  \node[below] at (7.6,0) {\footnotesize \val{netsize}};
  \node[align=center,font=\footnotesize] at (5,2.1)
    {readable and editable in part \hspace{4mm} $\longrightarrow$ \hspace{4mm}
     neither};
\end{tikzpicture}
\caption{Where the two paradigms sit on the description-length axis, at
comparable playing strength. \emph{Schematic; both values are placeholders.} The
comparison is not that one paradigm is smaller but that the difference is one of
scale, which is what makes description length the variable this paper measures.
The strength condition is load-bearing: a size comparison between policies of
different strength says nothing.}
\label{fig:size}
\end{figure}

但这种优势需要持续的人力投入。增加规则并不难，难的是不断检查它是否破坏已有行为；程序越复杂，维护所需的人力越多。瓶颈不是程序不能表达，而是维护需要人力，无法长期持续投入。

\begin{center}
\fcolorbox{brandaccent}{white}{\parbox{0.93\linewidth}{\small
\textbf{观察 2：维护需要人力。} 程序的持续修改依赖人力；当人力无法长期投入，策略改进就会停止。}}
\end{center}

多人对抗环境高度变化：对手策略不断切换，使每一种局面都只能获得有限样本。没有对手的单机游戏中，RL 面对近似不变的环境，反复采样相对容易；而在对手多变、样本稀疏的环境中，单纯拟合每种行为模式很难充分学习。此时，从每条轨迹中提取结构并直接优化程序，能够把一次观测中的信息迁移到后续决策，因而更适合少样本条件。

LLM-based coding agent 可以同时利用并弥补前述两点。程序本身是高度启发性且可解释的输入；agent 读取逐决策的轨迹回放来维护程序。轨迹、程序与 policy 修改构成一条信息传递链：若中间表示不能保留决策细节，信道容量就会很低；程序形式则把这条链条直接打开，使回放中的信息能够指向具体规则并落实为局部修改。与此同时，agent 可以持续修改同一个可执行、可读、可局部编辑的程序，逐步承担原本需要人力完成的维护工作。

\begin{center}
\fcolorbox{brandaccent}{white}{\parbox{0.93\linewidth}{\small
\textbf{观察 3：LLM agent 的闭环作用。} 直接读取程序提供高度启发性、可解释的输入；结合轨迹回放，agent 获得宽信息通道来维护程序，并逐步承担原本需要人力的持续维护。}}
\end{center}

% THE NAMING AND THE CONCESSION, moved here from 02-preliminary.tex on 2026-08-20 when that
% section was deleted. The paragraph above already argued the channel is wide; what only
% existed in Section 2 was the NAME and the limit on what the name licenses. Both have to
% survive, and this is the last paragraph that can carry them -- Section 3.3's "observation
% trace" and Section 4's replay contract both use the term as though it were defined.
%
% THE CONCESSION IS THE POINT, and it must not be softened into a hedge. The channel's width
% is a property of what is available and of the agent reading it, NOT of the policy being
% symbolic: echo2025hindsight rewrites failed episodes into a memory, holds no program
% anywhere, and gets the same kind of per-sample gain. So the paper states two bounds over
% two different objects and never merges them -- lower over what the loop can read, where
% the symbolic form earns no credit; upper over what the loop can edit, where the program is
% answerable. Merging them is what produces the unsupportable version of the claim, and a
% reader will supply the merged version unaided if this is left loose. 01-introduction.tex
% already concedes the same limit; keep the two statements consistent if either is edited.
这里的“宽信道”需要轨迹提供丰富的决策信息，也需要程序形式把这些信息完整传递到 policy 修改；没有程序、只把失败经验写入记忆的循环也可能获得部分类似的信息增益~\citep{echo2025hindsight}。本文真正要测量的是：在对抗式游戏的少样本条件下，这条信息链能否通过一个可解释、可持续修改的程序转化为策略性能。

\end{CJK*}

\subsection{Introducing Adversarial Heuristic Learning}
\label{sec:hl-method}

We use \emph{Adversarial Heuristic Learning} (\ahl) for the following setting:
an agent iteratively edits one executable policy, receives decision-level traces
from matches against changing external opponents, and is scored by a referee-side
evaluation pool it cannot rewrite. AHL is therefore the adversarial-game
instantiation studied here, rather than a definition of heuristic learning in
general. The policy artifact persists across rounds; each revision is made from
the current program, the game rules, and the feedback collected so far.

The instantiation is a loop between two agents, and it is deliberately simple. A
\emph{policy agent} is the program under revision; it plays the game and produces
interaction traces. A \emph{generative agent} observes those traces, together with
the rules and the policy agent's own source, and either generates a policy or
revises the one in hand. Its output is the next version of the policy agent, which
is then rescored against the frozen test pool, and the loop repeats.
Algorithm~\ref{alg:loop} states it in full.
The complete overview, including the contest-game suite in which the loop is measured,
is shown in Figure~\ref{fig:hl-arena-overview}.

% REPLACED the tikz block diagram with this algorithm block on 2026-08-19 at the
% user's direction. The diagram drew three boxes and three arrows, which is what the
% prose already says; the algorithm carries what the prose cannot -- the order of
% operations, where the budget is decremented, what is logged, and the fact that the
% generative agent's two modes are one call site with different inputs.
%
% Kept deliberately short, because the user's framing of the loop is "就这么简单".
% Resist adding the description-length protocol, the KL measurement, or the budget
% conversion here; those are separate protocols and belong in their own paragraphs
% later in 3.2. An algorithm that grows past one column stops being the thing that
% makes the loop look simple.
%
% [b] for the same reason as Figure 1: with [t] the float landed above its own
% subsection heading, inside the tail of 3.1. Check the rendered page after any edit
% above it.
\begin{algorithm}[b]
\caption{Heuristic learning. A generative agent observes the interaction between
the policy agent and its environment, then generates or revises the policy agent.
One revision per iteration; every version, score, and trace is appended to the log.}
\label{alg:loop}
\begin{algorithmic}[1]
  \Require game rules $R$, training pool $\mathcal{O}_{\text{train}}$, frozen test
    pool $\mathcal{O}_{\text{test}}$, budget $B$
  \State $\pi_0 \gets \Call{Generate}{R}$
    \Comment{no policy in hand: the agent writes one}
  \State $\mathcal{L} \gets \langle\,\rangle$
    \Comment{append-only log}
  \For{$t = 1$ \textbf{to} $B$}
    \State $\tau_t \gets \Call{Play}{\pi_{t-1}, \mathcal{O}_{\text{train}}}$
      \Comment{interaction traces, including the losses}
    \State $\pi_t \gets \Call{Revise}{\pi_{t-1}, R, \textsc{src}(\pi_{t-1}), \tau_t}$
      \Comment{reads; then edits}
    \State $e_t \gets \Call{Score}{\pi_t, \mathcal{O}_{\text{test}}}$
      \Comment{Elo on the frozen pool}
    \State append $(\pi_t, e_t, \tau_t)$ to $\mathcal{L}$
  \EndFor
  \State \Return $\mathcal{L}$
    \Comment{the trajectory is the result, not $\pi_B$ alone}
\end{algorithmic}
\end{algorithm}

We use one revision per iteration and append every policy version, score, and trace
to the log so that changes remain attributable.

\subsection{Key Elements of AHL}
\label{sec:hl-bottleneck}

Algorithm~\ref{alg:loop} has two elements that determine what AHL can learn. The
maintenance rule determines whether feedback can be written into the persistent
program continuously and effectively. The feedback information determines how much
the agent can observe in each round, including which opponents it plays and which
parts of each interaction are recorded.

\paragraph{Maintenance rule.}
The rule controls how revisions are proposed, checked, accepted, and remembered.
It must prevent a fix for one opponent from silently breaking behaviour that already
works against others, while retaining rejected changes so that the loop does not
rediscover them. Rescoring each revision against a fixed reference and logging the
decision makes the maintenance process attributable; its stability is an empirical
question evaluated in Section~\ref{sec:experiments}.

\paragraph{Feedback information.}
The feedback contract controls the input channel. Outcome-only feedback forces the
agent to guess which rule caused a loss, whereas decision-level traces expose the
sequence of actions, opponent replies, and the turning point needed to edit a
specific rule. In adversarial games, the opponent pool is part of this contract:
the agent can only learn from failures that the selected opponents elicit. Thus both
who the policy plays and what each replay contains determine the information
available for the next revision; our concrete trace is specified in
Section~\ref{sec:benchmark}.
